\documentclass[9pt,a4paper]{extarticle}

% ---------- Packages ----------
\usepackage[
  top=0.5cm,
  bottom=0.5cm,
  left=1cm,
  right=1cm
]{geometry}

\usepackage[utf8]{inputenc}
\usepackage[T1]{fontenc}
\usepackage{xcolor}
\usepackage{enumitem}
\usepackage{hyperref}
\usepackage{graphicx}
\usepackage{tikz}

% ---------- Couleurs ----------
\definecolor{navy}{HTML}{1F2A3C}
\definecolor{darkblue}{HTML}{0A1628}
\definecolor{gray}{HTML}{606060}
\definecolor{black2}{HTML}{1A1A1A}
\definecolor{accent}{HTML}{93A3B3}

% ---------- Réglages ----------
\pagestyle{empty}
\renewcommand{\familydefault}{\sfdefault}
\setlength{\parindent}{0pt}
\setlength{\parskip}{2pt}
\linespread{1}

\hypersetup{
  colorlinks=true,
  urlcolor=navy,
  linkcolor=navy
}

% ---------- Photo et nom ----------
\newcommand{\cvname}[2]{%
  \noindent
  \begin{minipage}[c]{3.1cm}
    \centering
    \begin{tikzpicture}
      \clip (0,0) circle (1.5cm);
      \node[inner sep=0pt] at (0,0) {%
        \includegraphics[
          width=3cm,
          height=3cm,
          keepaspectratio
        ]{ma-photo.jpeg}%
      };
    \end{tikzpicture}
  \end{minipage}%
  \hspace{0.3cm}%
  \begin{minipage}[c]{\dimexpr\linewidth-3.4cm\relax}
    \centering
    {\fontsize{28}{31}\selectfont
      \bfseries\color{navy}#1\par}
    \vspace{6pt}
    {\fontsize{8.5}{10.5}\selectfont
      \color{gray}#2\par}
  \end{minipage}
  \par\vspace{5pt}
  {\color{navy}\hrule height 0.8pt}
}

% ---------- Présentation agrandie ----------
\newcommand{\cvprofile}[1]{%
  \par\vspace{7pt}
  {\centering
    \fontsize{12}{14}\selectfont
    \itshape\color{black2}#1\par}
  \vspace{3pt}
}

% ---------- Rubriques ----------
% Répartit l'espace disponible entre les rubriques.
\newcommand{\cvsection}[1]{%
  \par
  \vspace{4pt plus 1fill}
  {\bfseries\large\color{navy}\MakeUppercase{#1}\par}
  \vspace{1pt}
  {\color{accent}\hrule height 0.5pt}
  \vspace{3pt}
}

% ---------- Titres des entrées ----------
\newcommand{\cventry}[3]{%
  \noindent
  {\small\bfseries\color{darkblue}#1}%
  \hfill
  {\footnotesize\itshape\bfseries\color{accent}#3}\par
}

% ---------- Sous-titres ----------
\newcommand{\cvsub}[1]{%
  {\footnotesize\itshape\color{gray}#1\par}
  \vspace{1pt}
}

% ---------- Listes ----------
\newenvironment{cvitemize}{%
  \begin{itemize}[
    leftmargin=10pt,
    itemsep=1pt,
    topsep=1pt,
    parsep=0pt,
    partopsep=0pt,
    label=\textbullet,
    font=\color{accent}
  ]
  \fontsize{8}{9.2}\selectfont
  \color{black2}
}{%
  \end{itemize}
  \vspace{2pt}
}

% ---------- Texte mis en valeur ----------
\newcommand{\hl}[1]{%
  \textbf{\color{darkblue}#1}%
}

\begin{document}

% Hauteur fixe : l'espace est réparti jusqu'en bas de la page.
\noindent
\begin{minipage}[t][\dimexpr\textheight-2pt\relax][s]{\textwidth}

% ================= EN-TETE =================

\cvname{Nisrine KHNIFASS}{%
  \href{mailto:nisrine.khnifass@hotmail.com}
       {nisrine.khnifass@hotmail.com}
  ~\textbar~
  +33 7 67 87 07 45
  \\[3pt]
  LinkedIn~: Nisrine Khnifass
  ~\textbar~
  Longvic, France
  \\[3pt]
  GitHub~:
  \href{https://github.com/NisrineKhnifass}{NisrineKhnifass}
  \\[3pt]
  Portfolio~:
  \href{https://nisrinekhnifass.github.io/portfolio/}
       {https://nisrinekhnifass.github.io/portfolio/}
}

\cvprofile{%
  Élève ingénieure en dernière année à l'ENSICAEN en
  développement informatique, spécialisée en cybersécurité
  et e-paiement, actuellement en échange académique au Japon.
  À la recherche d'un stage de 6 mois en cybersécurité
  à partir de mars 2027.
}

% ================= EXPERIENCE =================

\cvsection{Expérience --- Informatique \& Ingénierie}

\cventry
  {Stagiaire recherche --- Sungkyunkwan University (SKKU)}
  {}
  {Avril -- Août 2026}

\cvsub{Séoul, Corée du Sud}

\begin{cvitemize}
  \item \hl{Conçu et livré un prototype fonctionnel} de système
  de surveillance des dangers domestiques pour enfants,
  opérationnel de bout en bout à partir d'un flux vidéo
  de caméra de surveillance

  \item Pipeline en 5 modules (détection enfant,
  reconnaissance d'activité, contexte, décision du risque,
  notification) en \hl{Python (PyTorch/TensorFlow)}, avec un
  module de perception traitant les flux vidéo à
  \hl{\textasciitilde48 images/seconde}

  \item Atteint \hl{88,7\,\% d'accuracy} et
  \hl{100\,\% de rappel sur les situations dangereuses}
  (aucun danger manqué sur 335 vidéos de test)~;
  comparé 9 modèles de détection, 3 modèles de reconnaissance
  vidéo et 6 configurations de décision (dont 4 LLM)
\end{cvitemize}

% ================= PROJETS =================

\cvsection{Projets}

\cventry
  {Scanner de détection Magecart --- Projet personnel}
  {}
  {2026 -- en cours}

\begin{cvitemize}
  \item Développement d'un scanner automatisé des scripts tiers
  sur les pages de paiement e-commerce, en \hl{Python}.
  \hl{Résultat attendu~:} identification des cas de web-skimming
  avant exploitation

  \item Implémentation d'un module de conformité aux exigences
  \hl{PCI DSS 6.4.3 / 11.6.1} (inventaire et détection
  d'altération de scripts).
  \hl{Résultat attendu~:} rapport de conformité exploitable
\end{cvitemize}

\cventry
  {Cheffe de projet --- Escape Game interactif --- H3P Solution}
  {}
  {2026}

\begin{cvitemize}
  \item \hl{Géré et piloté une équipe de 4 personnes}
  via \hl{GitLab} (gestion de version, suivi des tâches)~:
  conception d'une application mobile \hl{Flutter},
  escape game interactif

  \item Développé des fonctionnalités avancées~:
  gestion simultanée de \hl{plusieurs équipes en temps réel},
  carte auto-générée, chronomètre, questions interactives

  \item \hl{Livré un MVP fonctionnel}, validé et apprécié
  par le client
\end{cvitemize}

\cventry{Sudoku Assistant --- Qt / C++}{}{2026}

\begin{cvitemize}
  \item Développé un assistant de résolution de grilles
  de Sudoku, avec une \hl{architecture MVP}

  \item Implémenté une internationalisation dynamique
  de l'interface (français / anglais)
\end{cvitemize}

% ================= FORMATION =================

\cvsection{Formation}

\cventry
  {Shibaura Institute of Technology}
  {}
  {Sept. 2026 -- Fév. 2027}

\cvsub{Tokyo, Japon ~\textbar~ Semestre d'échange académique}

\begin{cvitemize}
  \item Poursuite du cursus d'ingénieur en informatique
  dans le cadre d'une mobilité internationale
\end{cvitemize}

\cventry
  {ENSICAEN --- École nationale supérieure d'ingénieurs de Caen}
  {}
  {2024 -- 2027}

\cvsub{%
  Caen, France ~\textbar~
  Diplôme d'ingénieur en développement informatique,
  spécialisation Cybersécurité
}

\begin{cvitemize}
  \item Cours~: Cybersécurité, Cryptographie,
  Sécurité du logiciel, Technologie de sécurité,
  Théorie des codes, Monétique \& Paiement électronique
  (ISO7816, EMV, PCI DSS), Génie logiciel,
  Développement web \& mobile,
  Fouille de données \& apprentissage
\end{cvitemize}

\cventry{EM Normandie (double diplôme)}{}{2025 -- 2027}

\cvsub{Caen, France ~\textbar~ Master's Degree in Management}

\begin{cvitemize}
  \item Cours~: Leadership \& Management, Marketing,
  Stratégie d'entreprise, Finance, Éthique
\end{cvitemize}

\cventry
  {Classes préparatoires, Lycée Carnot --- mention Bien}
  {}
  {2022 -- 2024}

\cvsub{%
  Dijon, France ~\textbar~
  Préparation aux concours des grandes écoles d'ingénieurs
  (PCSI -- PC)
}

\begin{cvitemize}
  \item Algèbre, Analyse, Informatique,
  Physique fondamentale, Chimie organique et physique
\end{cvitemize}

% ================= COMPETENCES =================

\cvsection{Compétences}

{\fontsize{8.5}{10}\selectfont
\textbf{\color{navy}Soft Skills~:}
adaptabilité, esprit d’initiative, aisance à l’oral,
curiosité technique, coordination d’équipe\\

\textbf{\color{navy}Cybersécurité~:}
cryptographie, sécurité du logiciel,
technologie de sécurité, théorie des codes\\
\textbf{\color{navy}Monétique \& Paiement~:}
cartes à puce (ISO7816, EMV),
architecture \& écosystème du paiement électronique,
flux monétique, fraude \& PCI DSS\\
\textbf{\color{navy}Langages~:}
Python, Java, C, C++, JavaScript, Dart, Bash\\
\textbf{\color{navy}Frameworks \& bibliothèques~:}
Flutter, React, Node.js, Qt\\
\textbf{\color{navy}Bases de données~:}
PostgreSQL\\
\textbf{\color{navy}Data / IA~:}
PyTorch, TensorFlow, vision par ordinateur,
fouille de données et apprentissage\\
\textbf{\color{navy}Outils \& Méthodes~:}
Git, GitLab, GitHub, Linux ---
SOLID, patrons de conception, Agile\\

\textbf{\color{navy}Langues~:}
Français (natif), Anglais (courant, TOEIC 900),
Espagnol (intermédiaire), Allemand (débutant)\par
}

% ================= AUTRES EXPERIENCES =================

\cvsection{Autres expériences professionnelles}

\cventry
  {Vendeuse conseil --- C\&A}
  {}
  {Août 2024 -- Juillet 2025}

\begin{cvitemize}
  \item Gestion autonome d'un rayon~:
  réassort, contrôle des stocks,
  conseil et accompagnement client
\end{cvitemize}

\cventry
  {Coordinatrice de location immobilière --- Indépendante}
  {}
  {2022 -- 2024}

\begin{cvitemize}
  \item Gestion de locations d'appartements,
  suivi des contrats, planification et tâches administratives
\end{cvitemize}

% ================= INFORMATIONS COMPLEMENTAIRES =================

\cvsection{Informations complémentaires}

{\fontsize{8.5}{10}\selectfont
\textbf{\color{navy}Centres d'intérêt~:}
Entrepreneuriat, Voyage, Sport, Cuisine
~\textbar~
\textbf{\color{navy}Divers~:}
Permis B (véhiculée), BAFA\par
}

\end{minipage}

\end{document}